% !PRINTED 54
\weberpage{54}

3. If $u_0$ is a point of discontinuity, then there is a \gesperrt{positive integer} $m$ of such a kind that $(u-u_0)^m\varphi(u)$ remains finite, continuous and different from zero at the point $u_0$, and is therefore developable in integral positive powers of $u - u_0$. In this case the point of discontinuity is said to be \gesperrt{of the $m$th order}. From this there results a development of $\varphi(u)$ in ascending powers of $u - u_0$ which begins with $(u-u_0)^{-m}$ and converges within a circle described about $u_0$ that reaches as far as the nearest point of discontinuity. The coefficient of $(u-u_0)^{-1}$ in this development is called the \gesperrt{residue} of this function for the point $u_0$.

4. \gesperrt{Cauchy's theorem.} The integral
\[
\frac{1}{2\pi i}\int \varphi(u)\,du,
\]
extended in the positive sense over the boundary of a finite piece of surface which contains no points of discontinuity on its boundary line, is equal to the sum of the residues for the points of discontinuity lying in the interior of the piece of surface. The positive direction of integration is here to be taken as that one which lies with respect to the interior of the piece of surface as the $v$-axis lies to the $w$-axis, so that, with the usual mode of determination of these axes, the interior of the surface remains on the left when the boundary is traversed in the positive sense.

5. The function
\[
\psi(u) = \frac{d\log\varphi(u)}{du} = \frac{1}{\varphi(u)}\frac{d\varphi(u)}{du}
\]
has only points of discontinuity of the first order, and if at a point $u_0$ the product $(u-u_0)^{-m}\varphi(u)$ is finite and different from zero, then $m$ is the residue of $\psi(u)$ for this point. If $m$ is positive, then $u_0$ is called a zero point of the $m$th order of $\varphi(u)$.

Accordingly the following is an immediate consequence of \gesperrt{Cauchy's} theorem:

6. The integral
\[
\frac{1}{2\pi i}\int d\log\varphi(u),
\]
extended in the positive direction over the boundary of a piece of surface, is equal to the number of zero points, diminished by the number of points of discontinuity which lie in the interior of this
%CONT>>
